%%%%%%%%%%%%%%%%%%%%%%%%%%%%%%%%%%%%%%%%%%%%%%%%%%%%%%%%%%%%%%%%%%%%%%%%%%%%%%
%% Neural Computing and Applications (Springer) LaTeX template
%% Journal: Neural Computing and Applications
%% Reference style: Numbered (sn-mathphys-num)
%%%%%%%%%%%%%%%%%%%%%%%%%%%%%%%%%%%%%%%%%%%%%%%%%%%%%%%%%%%%%%%%%%%%%%%%%%%%%%

\documentclass[pdflatex,sn-mathphys-num]{sn-jnl}

% ── Packages ──────────────────────────────────────────────────────────────────
\usepackage{graphicx}
\usepackage{multirow}
\usepackage{amsmath,amssymb,amsfonts}
\usepackage{amsthm}
\usepackage{mathrsfs}
\usepackage[title]{appendix}
\usepackage{xcolor}
\usepackage{textcomp}
\usepackage{manyfoot}
\usepackage{booktabs}
\usepackage{algorithm}
\usepackage{algorithmicx}
\usepackage{algpseudocode}
\usepackage{listings}
\usepackage{array}
\usepackage{enumitem}
\usepackage{longtable}
\usepackage{pdflscape}
\usepackage[table]{xcolor}
\usepackage{url}

% ── Code listing style ────────────────────────────────────────────────────────
\lstset{
    basicstyle=\ttfamily\footnotesize,
    backgroundcolor=\color{gray!8},
    frame=single,
    rulecolor=\color{gray!40},
    breaklines=true,
    breakatwhitespace=true,
    numbers=none,
    xleftmargin=4pt,
    xrightmargin=4pt,
    aboveskip=4pt,
    belowskip=4pt,
    columns=fullflexible,
}

\theoremstyle{thmstyleone}
\newtheorem{theorem}{Theorem}
\newtheorem{proposition}[theorem]{Proposition}

\theoremstyle{thmstyletwo}
\newtheorem{example}{Example}
\newtheorem{remark}{Remark}

\theoremstyle{thmstylethree}
\newtheorem{definition}{Definition}

\raggedbottom

%%%%%%%%%%%%%%%%%%%%%%%%%%%%%%%%%%%%%%%%%%%%%%%%%%%%%%%%%%%%%%%%%%%%%%%%%%%%%%
\begin{document}

\title[BRIDGE: RAG Pipeline Design and Retrieval Evaluation] {Developing a Multi-Agent Retrieval-Augmented System for Interprofessional Healthcare Education Settings: The BRIDGE Platform}

%%=============================================================%%
%% Author block --- anonymize for double-blind submission.
%%=============================================================%%
\author*[1]{\fnm{Anusha} \sur{Umashankar}}\email{anushau@gwu.edu}
\equalcont{These authors contributed equally to this work.}

\author*[1]{\fnm{Hari Prasannaa} \sur{Thangavel Ravi}}\email{hariprasannaa.thangavelravi@gwu.edu}
\equalcont{These authors contributed equally to this work.}

\author[1]{\fnm{Amir} \sur{Jafari}}\email{ajafari@gwu.edu}

\author[2]{\fnm{Andrew} \sur{Wiss}}\email{awiss@gwu.edu}

\affil[1]{\orgdiv{Data Science Program}, \orgname{The George Washington University},
          \orgaddress{\city{Washington}, \state{DC}, \country{USA}}}

\affil[2]{\orgdiv{Milken Institute School of Public Health}, \orgname{The George Washington University},
          \orgaddress{\city{Washington}, \state{DC}, \country{USA}}}

%%=============================================================%%
\abstract{While interprofessional education (IPE) equips health care professionals with competencies in interprofessional collaboration, scalable team-based work proves challenging. In response to this need, we present the BRIDGE. BRIDGE (Bridging Roles and Responsibilities through Interdisciplinary Dialogue and aGent Engagement) consists of a multi-agent conversational system where each agent represents one of six professional roles within the healthcare system: Nurse, Physician Assistant, Medical Social Worker, Physical Therapist, Public Health Professional, and Health Administrator. These agents respond to a rich case study containing clinical, human and health systems perspectives based on their corresponding professional roles. To decrease the likelihood of hallucinations during interactions with BRIDGE agents, responses are grounded in a profession-specific, curated corpus via a Retrieval Augmented Generation (RAG) layer. We contribute here two elements: first, a novel dynamic prompt engineering approach incorporating iterative prompt engineering, runtime injection of RAG chunk, user profile personalization, and referral-based prompts; second, a comprehensive study exploring the impact on performance of the design decisions that have the greatest impact on the retrieval quality: document chunking, embedding model, retrieval depth, and metadata-based pre-filtering. We compare five different chunking techniques: fixed, fixed with overlap, recursive, sentence pack, and semantic, each with four different token targets: 3000, 5000, 7000, and 10000 tokens against four different retrieval depths ($k \in \{3,5,7,9\}$) in a corpus of 116 chunks extracted from 75 source documents. We evaluate seven different embedding approaches, consisting of five different dense neural models (MPNet-base-v2, BGE-small-en, MiniLM-L6-v2, Instructor-XL, and BGE-M3) and two baseline sparse lexical models (BM25, TF-IDF) using the best-performing retrieval depth k. Two different filters are also evaluated: topic-based filtering and profession-based filtering. Recursive chunking with $k{=}7$ a 7000 token target provided the best balance of hit rate (0.867) and prompt length. The dense neural embedding model BGE-M3 achieved the highest F1 (0.243), with BM25 providing perfect MRR. Profession-based filtering reduced the size of each individual agent’s corpus by an order of magnitude. Topic-based filtering did not significantly improve results. Discussion of trade-offs involved in making these decisions,
and next steps in validating this design approach are discussed.}

\keywords{retrieval-augmented generation, interprofessional education, embedding models, opioid use disorder, dense retrieval, BM25, multi-agent LLM, healthcare AI}

\maketitle

%%=============================================================%%
\section{Introduction}\label{sec:intro}

Healthcare outcomes for people with complicated medical cases depend on collaboration and communication between professionals from different domains. That is why Interprofessional Education (IPE) exists, it allows trainees for different health sector occupations, such as medicine, nursing, pharmacy, social work, and allied professions, to come together to learn how to collaborate even before joining the healthcare sector \cite{reeves2016ipe,who2010ipe}. The problem here is rarely lack of motivation but organization. Getting trainees from different schools to join in one activity is costly because of scheduling problems and costs of recruitment of simulated patients as well as hosting the discussion with a facilitator in a small group format.

LLMs offer a new solution. An artificial conversation partner will always be available, will have enough time to exchange information with a student, and can be customized through a detailed system prompt to impersonate certain professionals. Recent research has proven that several agents can be programmed to cooperate with each other while being instructed individually \cite{park2023generative,wu2023autogen}. In a sense, such a model would imitate what IPE tries to achieve: having six agents responding to the same query in parallel will create multi-perspectivity, which should serve as preparation for face-to-face interaction with professionals.

Two obstacles stand in the way. Health care sector information is unstable – opioid prescribing restrictions, overdose thresholds, and harm reduction advice change faster than models are trained. Second, large language models are prone to generating facts without a shadow of doubt, which is particularly hazardous where there is no room for error: medication dosage, risk of legal trouble, and emergency scenarios \cite{ji2023survey,singhal2023medpalm}. A model that tells the learner an erroneous dose of naloxone cannot be considered an educational tool.

A Retrieval-Augmented Generation (RAG) paradigm can solve both problems: rather than relying on the information learned during the training phase, the model will search a manually curated collection of texts and use evidence from that collection to answer the query \cite{lewis2020rag,gao2023ragsurvey}. The quality of a particular RAG model depends on multiple choices made about the process: document chunking, selection of an embedding model for ranking chunks, number of retrieved chunks, and metadata-based retrieval. Such a process requires more rigorous analysis in the case of a highly sensitive domain like clinical education.

In this paper, we present our research into creating RAG for the BRIDGE project, a multi-agent conversational model for interprofessional education on opioid use disorder. Specifically, we (i) propose a per-agent retrieval architecture with individual retrieval of a separate part of the corpus for each agent; (ii) create a dynamic prompt generation pipeline addressing shortcomings of static prompts: iterative prompt engineering, referral awareness, and personalized user profiles; (iii) benchmark five chunking algorithms on four token budgets and four retrieval depth settings; (iv) test seven retrieval techniques, including both sparse lexical retrieval methods and dense embeddings, on the opioid corpus; and (v) assess two approaches to metadata-based retrieval in the BRIDGE context. We do not attempt to propose a new retrieval algorithm but rather provide a reproducible comparison for educational tool developers in similar domains. Measuring downstream learning outcomes is outside the scope of this work and is reserved for a planned user study.

The rest of the paper is structured as follows. Section~\ref{sec:related} gives an overview of prior work in RAG, multi-agent LLMs, and computational support for IPE. Section~\ref{sec:bridge}  outlines the baseline BRIDGE system and describes its drawbacks. Section~\ref{sec:proposed} introduces the proposed per-agent retrieval architecture. Section~\ref{sec:methods} covers corpus, extraction pipeline, and experiment design. Section~\ref{sec:results} provides the results of chunking algorithm comparisons, retrieval performance benchmarks, and metadata filter testing. Section~\ref{sec:implementation} discusses implementation details and referral mechanisms. Section~\ref{sec:discussion} discusses results and limitations, and Section~\ref{sec:conclusion} concludes.

%%=============================================================%%
\section{Related Work}\label{sec:related}

\subsection{Retrieval-Augmented Generation}

RAG aimed at unifying the parametric memory in a trained language model with a non-parametric store where it could reference up-to-date evidence instead of recalling memorized knowledge from training \cite{lewis2020rag}. Izacard and Grave followed up with Fusion-in-Decoder, where the retrieved passages were encoded separately from each other before being decoded \cite{izacard2021fid} Gao et al. conducted a survey, describing the development of the field in three stages from naive to advanced and modular approaches, where the latest studies concentrated on better retrieval accuracy, multi-hop search capabilities, and minimization of prompts \cite{gao2023ragsurvey}.

The chunking strategy is not getting enough consideration in research literature compared to its importance. Practice-oriented papers describe how practitioners have agreed on a handful of heuristics for implementing this step: window-size-based segmentation, recursive character partitioning, and semantically aware embedding segmentation, all of which are available off-the-shelf in frameworks like LangChain and LlamaIndex \cite{langchain2023,kamradt2023}. Late chunking techniques and small-to-big retrieval algorithms demonstrate how this step can be researched independently from retrieval in a few modern papers \cite{gunther2024late}. In our study, we pragmatically assume that there are several chunkers already known to work well and evaluate their performance on clinical data.

\subsection{Embedding Models for Dense Retrieval}

Similarity scores for dense retrievers are computed via cosine similarity of learned vector representations of documents. The Sentence-BERT \cite{reimers2019sbert} and MPNet \cite{song2020mpnet} models proved the feasibility of transformers in sentence encoding. Later models such as BGE series \cite{xiao2023bge,chen2024bgem3} and E5 \cite{wang2022e5} made improvements in terms of retrieval perfection on benchmarks from general domains. Sparse lexical approaches still show good results. In cases where there is a need for term overlap, BM25 proves to be a robust baseline approach \cite{robertson2009bm25}. Recent research warns against blindly believing in dense retriever performance without proof \cite{thakur2021beir}.

\subsection{Multi-Agent LLM Systems}

Coordinating multiple agents driven by LLMs to take on different responsibilities remains a challenging research topic. Generative agents, for example, have been found to display realistic human behavior in a simulated social environment \cite{park2023generative}. Formal systems like AutoGen \cite{wu2023autogen} and CAMEL \cite{li2023camel} make use of a dialogue between multiple agents when dealing with program development and other problem-solving exercises. What sets the IPE application apart from others is that its objective is not achieving efficient agreement but identifying role-based reasoning to be contrasted by humans.

\subsection{AI for Healthcare Education}

MedPaLM \cite{singhal2023medpalm}, BioBERT \cite{lee2020biobert}, or any other medical LLM has achieved excellent results in medical question-answering evaluation tasks. The evidence from education research is less abundant. Virtual standardized patients can be traced back even before the emergence of LLMs \cite{stevens2006virtual}, and some more recent work has looked at LLM-based intelligent tutors in clinical scenarios. To our best knowledge, there exists no prior work that implements such a profession-partitioned per-agent RAG architecture for IPE, and there is no analysis of such trade-offs for an opioid education dataset. Relevant
papers addressing related issues of interest to this work in the area of IPE, AI acceptability in collaborative practice, and use of AI in the collaboration process include:

\begin{itemize}
    \item Wiss et al. \cite{wiss2025groupthink} look into the utilization of AI to create productive controversies in collaborative problem-based learning environments.
    \item Wiss et al. \cite{wiss2025aaipe} created the AI Acceptance Scale for IPE and collaboration setting.
\end{itemize}

%%=============================================================%%
\section{The BRIDGE System}\label{sec:bridge}

The BRIDGE tool is a platform utilizing artificial intelligence and multiple agents designed to facilitate the process of IPE within the domain of healthcare education. The BRIDGE simulation is based on the collaboration among health professionals through the representation of six distinct agents corresponding to real-life professions: Nurse, Physician Assistant, Medical Social Worker, Physical Therapist, Public Health Professional, and Health Administrator.

\subsection{Architecture}

The backend consists of an entirely asynchronous FastAPI application that is connected to OpenAI’s GPT-5.2 via the \texttt{AsyncOpenAI} interface. The orchestration layer for multiple agents, called \textit{Interprofessional AgentSystem}, calls all the agents concurrently with the help of Python’s \texttt{asyncio.gather()} such that each perspective by the respective professional does not have to wait for the previous one to finish execution. Each agent retains its conversational consistency with the help of response-ID functionality provided by OpenAI. A common conversation memory is maintained in PostgreSQL, which, if it grows larger than 6{,}000 characters, is compressed asynchronously.

\subsection{Personalization}

During the registration process, the users fill out an initial questionnaire on their professional qualifications such as their current job position, years of work experience, educational background, and previous clinical involvement with opioids. The completed form is then placed in front of each agent’s prompt for generating appropriate answers.

\subsection{Limitations of the Baseline}

Despite its effectiveness, there are four main disadvantages that prompted this research:

\begin{enumerate}[leftmargin=2em]
\item \textit{Absence of external grounding.} All answers are provided based on the parameters set for GPT-4o-mini. Guidelines and protocols from hospitals are not used during the inference process, putting the system at risk of hallucinations and drifts in response due to the changes made in guidelines over time.
\item \textit{Limited scope of cases.} There is only one case (Julia Baker) presented to students, which does not allow for generalization of results. Multimorbid cases are not considered, and other specializations have not been studied yet.
\item \textit{Lack of adaptability of prompts.} System prompts are written by developers manually. There are no tools within the current pipeline to change prompts depending on users' performance or conversation dynamics.
\item \textit{No quantitative assessment.} Response quality, clinical accuracy, and learning effectiveness are not assessed. This is a major hurdle in validating BRIDGE as a training tool.
\end{enumerate}

In this paper, three out of four issues are tackled. First, the RAG component discussed in Section~\ref{sec:proposed} introduces grounding agents' responses using a specialized library for each field. Second, the patient case management interface (Section~\ref{sec:implementation}) allows instructors to include new patient scenarios apart from the first example of an opioid user. Third, there are four mechanisms addressing static prompts: iterative prompt tuning, where each agent's personality and communication style are developed iteratively; dynamic prompt generation, where selected RAG documents and users' queries are injected into the prompt template at inference time; user-specific prompts, where the users' professional profile and experience levels are added to the prompt of each agent; and referral prompts, where each agent recognizes a question outside their specialization and guides the user to consult another doctor.

\begin{figure}[h]
    \centering
    \includegraphics[width=0.85\linewidth]{currentarchi.pdf}
    \caption{BRIDGE baseline system architecture. A student query is authenticated, enriched with the user profile, and dispatched to the selected agents in parallel. Each agent calls GPT-4o-mini using only parametric knowledge, and responses are logged to PostgreSQL.}
    \label{fig:baseline}
\end{figure}

%%=============================================================%%
\section{Proposed Approach}\label{sec:proposed}

\subsection{Design Goals}

There are three main objectives of the RAG system that have been taken into account due to their relevance to the shortcomings observed in the baseline. The responses provided must be based on clinical evidence obtained from SAMHSA, CDC, FDA, and NIDA; the knowledge base must be separated such that each agent accesses content that pertains to its own occupation; and the corpus must be expandable without requiring retraining.

\subsection{System Overview}

BRIDGE is a web app built from end to end on a cloud platform (Table~\ref{tab:stack}). The frontend, written in React and TypeScript, offers the interface for the case study where students write the clinical situation and choose which agents to involve. The backend coded in Python FastAPI controls the calling of LLMs, handles conversation history, and passes retrieved evidence pieces to the corresponding agents.

\rowcolors{2}{gray!15}{white}
\begin{table}[h]
    \centering
    \caption{BRIDGE Technology Stack}
    \label{tab:stack}
    \begin{tabular}{@{}ll@{}}
        \toprule
        \textbf{Layer} & \textbf{Technology} \\
        \midrule
        Frontend        & React, TypeScript, Tailwind CSS \\
        Backend API     & FastAPI (Python) \\
        Database        & PostgreSQL (Users, ChatSessions tables) \\
        Authentication  & Supabase \\
        Hosting         & AWS EC2 + API Gateway + S3 \\
        LLM             & OpenAI \texttt{gpt-5.2} \\
        \botrule
    \end{tabular}
\end{table}

The novelty of this work falls in between the prompt creation step and the calling of LLMs, where retrieved pieces of evidence are injected into each agent’s context just before LLMs are called.

\subsection{Professional Agents}

The six agents cannot be substituted for each other; rather, each agent has a system
prompt that sets the tone for their professional background and mode of operation (Table~\ref{tab:agents}). It is immediately evident that, when compared side by side, there are striking differences between the approaches, such as the nurse’s concentration on airways,
the social worker’s concern regarding custody and continuity of care, and the public health agent’s concern with the population-based impact. Overall, the agents create the multiprofessional way of thinking that IPE aims to promote.

\rowcolors{2}{gray!15}{white}
\begin{table}[h]
    \centering
    \caption{BRIDGE professional agents and clinical focus.}
    \label{tab:agents}
    \begin{tabular}{@{}ll@{}}
        \toprule
        \textbf{Agent} & \textbf{Clinical Focus} \\
        \midrule
        Nurse                      & Patient care, nursing observations \\
        Physician Assistant        & Diagnostics, treatment planning \\
        Medical Social Worker      & Social care, community resources \\
        Physical Therapist         & Rehabilitation, recovery plans \\
        Public Health Professional & Disease prevention, health education \\
        Health Administrator       & Resource coordination, policy \\
        \botrule
    \end{tabular}
\end{table}

The six professions correspond to typical professions used within IPE models \cite{reeves2016ipe} and cover the continuum of care related to opioids.

\subsection{Per-Agent RAG Architecture}

Instead of relying on a single retriever that serves all the agents, the retrieval function is integrated within the execution of each individual agent. In this way, when the student fires off a question, it gets fired off simultaneously by the backend to all of the agents chosen by the user. Then each agent executes their own profession-filtered similarity search using the chunks associated with the chosen profession label.

\begin{figure}[h]
    \centering
    \includegraphics[width=0.85\linewidth]{ragflow.pdf}
    \caption{Proposed RAG architecture. A student query is routed in parallel to the selected profession agents. Each agent filters the corpus by its profession and topic tags, performs a dense similarity search, injects the top-$k$ passages into the prompt, and calls GPT-5.2.}
    \label{fig:rag_arch}
\end{figure}

This process leads to the generation of 5-6 evidence-backed answers to a student question, where each answer corresponds to a different professional perspective of the corpus. The corpus embeddings are pre-computed once at startup and stored in memory. This means that the actual similarity search takes less than one millisecond irrespective of the value of $k$.

%%=============================================================%%
\section{Methods}\label{sec:methods}

In this section, the dataset, the pipeline for extracting information from various source formats,
the chunking and embedding methods that were considered, and the criteria used to compare them are described.

\subsection{Corpus}

This corpus covers the clinical domain covered by all six BRIDGE interventions. At present, the focus is on Opioid Use Disorder (OUD). The corpus comprises 116 units taken from 75 documents, including 40 PDFs and 35 web pages. The PDF-based documents include federal clinical guidelines from SAMHSA, CDC, WHO, and CMS, WHO operational guidelines for the region, state policy documents, articles from PubMed Central and NIDA, reports related to CCBHC programs, and training manuals developed by an interprofessional team. The web-sourced portion was drawn from the CDC, NIDA, FDA, CHCS, NASHP, PMC, and the National Council for Mental Wellbeing.

\subsubsection{Profession Categorization}

A profession label or multiple profession labels are associated with each individual source document. A single URL that appears under more than one profession classification is accessed only once, and appropriate labels are associated with it. The chunk distribution for each profession is summarized in Table~\ref{tab:chunks_by_profession} below. Public Health, Physician Assistant, and Nurse are dominant professions due to the nature of the opioid literature, but Physical Therapist and Health Administrator are intentionally smaller
professions that the expansion work in Section~\ref{sec:discussion} aims to broaden.

\rowcolors{2}{gray!15}{white}
\begin{table}[h]
    \centering
    \caption{Approximate chunk distribution by profession.}
    \label{tab:chunks_by_profession}
    \begin{tabular}{@{}lc@{}}
        \toprule
        \textbf{Profession} & \textbf{Approximate chunks} \\
        \midrule
        Public Health Professional & $\sim$48 \\
        Physician Assistant     & $\sim$41 \\
        Nurse                   & $\sim$38 \\
        Medical Social Worker   & $\sim$21 \\
        Physical Therapist      & $\sim$6 \\
        Health Administrator    & $\sim$5 \\
        General                 & $\sim$30 \\
        \botrule
    \end{tabular}
    \footnotetext{A single chunk may carry multiple profession labels; column totals therefore exceed the corpus size of 116.}
\end{table}

\subsubsection{Opioid Topic Taxonomy}

After chunking, each chunk is automatically tagged with up to ten opioid-related topics using keyword matching (Table~\ref{tab:topics}). The taxonomy is intentionally coarse; a single chunk often carries several labels.

\rowcolors{2}{gray!15}{white}
\begin{table}[h]
    \centering
    \caption{Opioid topic taxonomy and representative keywords.}
    \label{tab:topics}
    \begin{tabular}{@{}ll@{}}
        \toprule
        \textbf{Topic} & \textbf{Representative keywords} \\
        \midrule
        \texttt{overdose}           & overdose, unresponsive, unconscious \\
        \texttt{emergency}          & call 911, life-threatening, emergency \\
        \texttt{naloxone}           & naloxone, narcan, intranasal \\
        \texttt{withdrawal}         & withdrawal, detox, taper, physical dependence. \\
        \texttt{dosage}             & dosage, mg, milligram, titrate, twice daily \\
        \texttt{treatment}          & buprenorphine, methadone, MOUD, OUD \\
        \texttt{prevention}         & harm reduction, prevention, safe storage \\
        \texttt{mental\_health}     & mental health, depression, PTSD \\
        \texttt{legal}              & controlled substance, schedule, HIPAA \\
        \texttt{patient\_education} & patient education, caregiver, warning signs \\
        \botrule
    \end{tabular}
\end{table}

\subsection{PDF Extraction}\label{sec:pdf_extraction}

The vast majority of the guideline documents exist in PDF format, which is well-optimized for presentation and not for structured data. Extracting these documents into clean, machine-read text was another non-trivial preprocessing step. We tested several extractors. PyPDF is fast but lacks proper structure extraction of tables and multi-columns resulting in disordered text where the reading order is incorrect quite often. Pdfplumber offers better extraction of table elements when the table borders are clear, but
falls behind on tables without borders or implied by spacing only. PyMuPDF offers better general-purpose text extraction thanks to using position information in text reconstruction, making it especially useful on multi-column documents. An additional OCR step using OCRmyPDF has been suggested for scanned or mixed PDF documents.

Docling was chosen as an extraction tool. The document understanding approach of Docling extracts more reliable structure in terms of tables, formatting, and reading order as compared to conventional tools at a non-trivial computational cost, which is justified due to one-time-only extraction for building each corpus.

\subsection{Web Extraction}

Every URL is downloaded using retry, timeout, and validation steps, and the source HTML content is cleaned using Trafilatura to get only the core content. After extraction, the content is converted into headings, paragraphs, lists, and tables. This process yields plain text that will be fed to chunking and embedding.

\subsection{Chunking Strategies}

Five chunking techniques were considered, namely, \textit{fixed-size}, \textit{fixed-size with overlap}, \textit{recursive}, \textit{sentence pack}, and \textit{semantic chunking}. For each of them, experiments were conducted with four target chunk sizes: 3000, 5000, 7000, and 10000 tokens. All techniques were tested with four retrieval depths ($k \in \{3, 5, 7, 9\}$) for a total of 80 combinations. As in the literature on RAGs \cite{langchain2023,kamradt2023,gunther2024late}, several measures were used: Recall@k (Hit rate), Average Rank of first hit, latency, chunk-build time, index-build time, and total chunk count.

\textbf{Fixed}: The text is segmented based on absolute word count, without any overlap.
\textbf{Overlap}: Like fixed segmentation, except subsequent segments have a 20\% overlap window between them. \textbf{Recursive}: The text is broken up according to paragraph borders first; for larger paragraphs, it is then further segmented according to sentence borders. \textbf{Sentence-Packing}: Sentences are greedily added to a single
chunk until the token budget is exceeded, after which point a new segment is started. \textbf{Semantic}: Sentence embeddings are calculated, and a new segment is
created whenever cosine similarity between sentences falls below a certain
threshold.

\subsection{Embedding Models}

Benchmarking was done using seven retrievers (Table~\ref{tab:models}). Dense neural encoders as well as sparse lexical encoders were included in the comparison. The dense neural encoders used were open-weight encoders that did not require any fine-tuning. The times listed in Table~\ref{tab:models} were for a one-time startup cost on the entire 116-chunk dataset.

\rowcolors{2}{gray!15}{white}
\begin{table}[h]
    \centering
    \caption{Embedding models: specifications and one-time encode time for the 116-chunk corpus.}
    \label{tab:models}
    \begin{tabular}{@{}llcc@{}}
        \toprule
        \textbf{Model} & \textbf{Dim.} & \textbf{Params.} & \textbf{Encode time} \\
        \midrule
        \texttt{all-mpnet-base-v2}    & 768   & 109M  & 29.6 s \\
        \texttt{BAAI/bge-small-en}    & 384   & 33M   & 20.2 s \\
        \texttt{all-MiniLM-L6-v2}     & 384   & 22M   & 4.2 s  \\
        \texttt{hkunlp/instructor-xl} & 768   & 335M  & 120 s  \\
        \texttt{BAAI/bge-m3}          & 1024  & 569M  & 150 s  \\
        BM25                          & ---   & ---   & 0.3 s  \\
        TF-IDF                        & vocab & ---   & 0.5 s  \\
        \botrule
    \end{tabular}
\end{table}

\subsection{Metadata Filtering}

Topic-based filtering and profession-based filtering approaches were assessed. The former approach filters the potential chunks based on their opioid topic tags that are aligned with the dominant topic in the query. In contrast, the latter limits the search for each agent to the manually curated chunks according to their professions (categories field).

\subsection{Evaluation Metrics}

For the embeddings comparison task, we use Mean Reciprocal Rank (MRR), Precision@$k$, Recall@$k$, and F1@$k$ metrics on a held-out evaluation set comprising representative student queries. MRR is considered the main metric since the retrieved passages are inserted verbatim into the prompt of the agent; hence, a high-quality top passage plays a bigger role compared to recall. For chunking comparison, Hit Rate and First Hit Rank are the main metrics.

Our evaluation query set was generated by the authors together with collaborators reviewing the work and includes questions from all ten categories described in the opioid taxonomy of Table~\ref{tab:topics}: overdose response, emergency procedures, administering naloxone, managing withdrawal, dose information, treatment modalities, harm reduction practices, mental issues, legal issues, and patient education. Questions are formulated in a way that is realistic to what healthcare students would ask the BRIDGE agents, encompassing both factual look-ups and scenarios requiring reasoning. We note that our set is purposefully limited to allow for a reproducible engineering analysis of different retrieval setups rather than statistical generalization.

%%=============================================================%%
\section{Experimental Results}\label{sec:results}

\subsection{Chunking and Retrieval Depth}

The average and the highest possible hit rate across all chunking configurations were 0.758 and 0.933, respectively. It seems that retrieval is reliable at moderately high k in this corpus, as the first relevant chunk is present on average in the top two retrievals. More detailed per-configuration statistics are provided in Appendix~\ref{appendix:chunking}.

An increase in $k$ value from 3 to 7 consistently improved hit rate. The difference in hit rate between $k=7$ and $k=9$ for most of the configurations was minor, as the well-known effect of diminishing returns means there is little useful information left for higher k beyond the first five or seven hits. For each chunking approach, there was only marginal improvement from adding $k$ after a certain point in terms of prompt overhead vs. coverage.

As opposed to an incremental increase in k, chunk size showed a significant effect on hit rate, although it was complicated by the chunking type. Chunking by semantic boundaries was effective for the largest chunk size used (10000 tokens), achieving a hit rate of 0.933 $k=5$ but yielding a significantly higher chunk count (61 chunks) and build time (60.7~s). On the other hand, fixed overlap and recursive chunking both achieved hit rate 0.933 for $k=9$ when the target chunk size was 10000 tokens but took less than 0.2~s to build a prompt.

The summary of our results in Table~\ref{tab:chunk_tradeoffs} shows the trade-off involved in choosing the optimal chunking configuration. We decided that the most efficient balance of retrieval quality and additional prompt cost is achieved using \textit{recursive chunking with target size 7000 tokens and $k=7$}. This approach yielded a hit rate of 0.867 for just 35 chunks and 0.149~s of build time. We will use this chunking approach in the comparison with embeddings below. It must be noted that this configuration might not prove optimal in larger corpora.

\rowcolors{2}{gray!15}{white}
\begin{table}[h]
\caption{Qualitative chunk-size trade-offs.}
\label{tab:chunk_tradeoffs}
\centering
\begin{tabular}{@{}p{4cm}p{4cm}@{}}
\toprule
\textbf{Smaller chunks} & \textbf{Larger chunks} \\
\midrule
More precise retrieval & Broader contextual coverage \\
More chunks to index & Fewer chunks to index \\
Higher indexing cost per corpus & Faster indexing \\
Risk of context fragmentation & Better context retention \\
\botrule
\end{tabular}
\end{table}

\subsection{Embedding Model Comparison}

The performance of all seven retrievers is presented at $k = 7$ on the unfiltered by professions corpus (Table~\ref{tab:results_k7}).

\rowcolors{2}{gray!15}{white}
\begin{table}[h]
  \centering
  \caption{Retrieval metrics at $k=7$. Best value per column in bold.}
  \label{tab:results_k7}
  \begin{tabular}{@{}lcccc@{}}
    \toprule
    \textbf{Model} & \textbf{MRR} & \textbf{Precision} & \textbf{Recall} & \textbf{F1} \\
    \midrule
    MiniLM-L6-v2   & 0.8000          & 0.6381          & 0.0900          & 0.1550          \\
    MPNet-base-v2  & 0.8000          & 0.6952          & 0.1220          & 0.1844          \\
    BGE-small-en   & 0.9333          & 0.7143          & 0.1020          & 0.1750          \\
    TF-IDF         & 0.9167          & 0.7619          & 0.1366          & 0.2076          \\
    BM25           & \textbf{1.0000} & \textbf{0.8286} & \textbf{0.1721} & 0.2426          \\
    BGE-M3         & \textbf{1.0000} & 0.8100          & 0.1630          & \textbf{0.2281} \\
    Instructor-XL  & 0.9333          & 0.6571          & 0.1015          & 0.1723          \\
    \botrule
  \end{tabular}
\end{table}

There are two clear observations. First, BM25 and BGE-M3 performed equally well on MRR = 1.000, indicating that all queries received relevant chunks at rank 1. Second, BM25 and BGE-M3 were also leaders in F1 scores (0.2426 and 0.2281, respectively); however, the difference from other retrievers is minor. TF-IDF achieved F1~=~0.208, which makes it a plausible option.

However, further consideration is required. While MRR = 1.000 is a ceiling effect on the 116-chunk dataset and cannot tell much about differences between retrievers, precision and recall at k = 7 can help make better distinctions. BM25 leads BGE-M3 on precision (0.829 vs. 0.810) and recall (0.172 vs. 0.163) at k = 7, though the difference is in the margin that is unlikely to be significant in practice. In general, dense retrievers tend to perform better than others on a larger dataset (the next step of our roadmap), which means that dense retrievers are preferable to other models because they demonstrate higher stability to vocabulary mismatch problems \cite{thakur2021beir,chen2024bgem3}. Therefore, BGE-M3 should be the default retriever for use cases requiring high-quality responses; otherwise, BM25 is a lightweight fallback model.

Nevertheless, the encoding costs must be taken into account. It takes 150~seconds to encode a 116-chunk dataset using BGE-M3. With respect to BM25, this number is 0.3~seconds.

\subsection{Metadata Filtering}

\subsubsection{Topic Filtering}

There was almost no gain in F1 scores using topic-based pre-filtering (Table~\ref{tab:topic_filter}). For BGE-M3, while F1 increased from 0.228 (no filtering) to 0.231 (topic filtering), such improvement cannot be considered significant. Topic filtering lowered F1 scores for BM25, TF-IDF, and Instructor-XL models from 0.004 to 0.008. Only MiniLM-L6-v2 and BGE-M3 had a marginal increase.

\rowcolors{2}{gray!15}{white}
\begin{table}[h]
  \centering
  \caption{Topic-filtered vs.\ baseline F1 at $k=7$.}
  \label{tab:topic_filter}
  \begin{tabular}{@{}lccc@{}}
    \toprule
    \textbf{Model} & \textbf{Baseline F1} & \textbf{Topic-filtered F1} & $\Delta$ \\
    \midrule
    MPNet-base-v2 & 0.1844 & 0.1844 & \phantom{$-$}0.0000 \\
    BGE-small-en  & 0.1750 & 0.1742 & $-$0.0008 \\
    BM25          & 0.2426 & 0.2385 & $-$0.0041 \\
    MiniLM-L6-v2  & 0.1550 & 0.1590 & $+$0.0040 \\
    TF-IDF        & 0.2076 & 0.1995 & $-$0.0081 \\
    Instructor-XL & 0.1723 & 0.1682 & $-$0.0041 \\
    BGE-M3        & 0.2281 & 0.2314 & $+$0.0033 \\
    \botrule
  \end{tabular}
\end{table}

We believe that the problem lies with the coarseness of our opioid taxonomy. Keyword-based chunking results in many chunks being assigned to more than one category. Sometimes, the query also maps to a category that is different from what we expect a person to categorize. Thus, single category pre-filtering does not capture the correct matches.

\subsubsection{Profession Filtering}

Profession filtering refers to a structural rather than a semantic constraint; each agent only considers those chunks curated for the discipline to which it belongs. For a dataset with 116 chunks, each agent retrieves roughly 30-40 chunks instead of all of them. In light of the above reasoning about how BRIDGE operates in production to serve each agent separately according to its evidence base and not the union, it becomes clear why this mode of filtering was chosen.

We have to make one observation. Since we compute our metrics of retrieval performance without profession filtering, they are reflective of the worst case. We expect the application of profession filtering to increase the value of recall at each particular $k$ for each agent (as the denominator decreases) but not affect MRR for the top performers for whom MRR equals 1.00 anyway. Evaluation of the system using ground-truth relevance labels per profession remains future work.

%%=============================================================%%
\section{System Implementation}\label{sec:implementation}

The components discussed in this section include three components that, while not technically retrieval experiments, are required to make BRIDGE an educational resource: the Retrieval and Generation Pipeline, Structured Referrals, and the Deployment Infrastructure.

\subsection{Retrieval and Prompt Generation Pipeline}

The pipeline is of a one-shot nature. With respect to a specific input query, the retriever fetches appropriate chunks using the input query string, the professional filter of the agent, and other information regarding the scenarios involved. These chunks get incorporated in the prompt of the agent along with the input query, after which GPT-5.2 gets invoked only once to generate the final output. Earlier, there was an attempt made at using two-shot architecture (invoking prompt-synthesis API and then the generation API separately), which proved unnecessary upon removing it.

\subsection{Agent Referral}

A professional who does not know the answer refers the patient to a professional who knows. Replicating this process on the agents means teaching the agents to identify when a question lies outside their expertise. Earlier versions detected this based on phrases like “you should ask the pharmacist” extracted by parsing the free-form reply. These failed to be reliably consistent: the smaller language model may skip the referral phrase completely or say something similar enough to escape detection. The current system uses OpenAI’s output JSON schema feature with \texttt{strict: true}:

\begin{lstlisting}
{ "response": "...", "refer_to": [...] }
\end{lstlisting}

The \texttt{refer\_to} field is an array whose elements are restricted to one or more of the constants: \texttt{nurse}, \texttt{physician}, \texttt{socialworker}, \texttt{therapist}, \texttt{publichealth}, \texttt{healthadmin}. Requiring this schema at the API level ensures that every single response is machine-parsable without worrying about changes due to models used. An agent is trained to generate a referral if the appropriate expert will provide a significantly more relevant response. The format of this referral consists of a brief statement followed by forwarding the same question to another agent.

When \texttt{refer\_to} is non-empty, the user interface provides links next to the response that refer to the mentioned professions. Upon clicking a link, the frontend sends back the initial question to that particular agent. To avoid circular referrals, the interface tracks all agents that have answered the question and disables buttons pointing to agents in the set of answers already received.

\subsection{Case Study Management}

A case management system that is administered by an instructor and can be accessed using the sidebar gives the ability to create new cases via a form where one is required to fill in the descriptive title, demographic information, a summary of the condition, and the clinical story that encompasses the present medical history, physical findings, lab results, and background information. Once added, a case is automatically added to the list of cases that can be selected by any user. Optional audio or video attachments (recorded patient interviews, procedure demonstrations) support richer, less text-bound scenarios.

\subsection{Infrastructure and Deployment}

The system uses AWS infrastructure. Backend services reside on Amazon EC2 instances and operate using the FastAPI orchestration layer. Processed document fragments and structured retrieval results are stored in Amazon RDS databases, whereas Supabase stores user records, authentication information, sessions, and application metadata. Asynchronous processing techniques are employed wherever possible to ensure that lengthy multi-agent processes do not impede the front-end interface. The division of the front-end and back-end layers enables scalability of both tiers independently as the corpus size and number of users increase.

%%=============================================================%%
\section{Discussion}\label{sec:discussion}

\subsection{Interpretation of Results}

Three conclusions can be drawn. Firstly, for a small-sized curated dataset, the best retrievers reach the saturation point of MRR. Users are cautioned not to overinterpret MRR equaling 1.000 to mean that one retriever is necessarily better than the other; saturation is simply a characteristic of the dataset size and the evaluation set used, rather than the model itself. More insight can be gained from precision, recall, and latency.

Secondly, sparse retrieval using classical BM25 still competes well with dense encoding models on this small-sized clinical dataset, owing to its relatively high vocabulary overlap between queries and passages due to the exactness of clinical language, including drug names, dosages, and ICD codes. Therefore, users are cautioned against assuming that dense retrievers always perform better than sparse ones automatically.

Lastly, the practice of filtering based on metadata can be categorized into two groups. Filtering based on semantic information such as topic tags obtained through keyword matching will likely remove valid documents from the corpus, and users are advised against this approach. Filtering based on structural information such as profession, which involves manual curation of the metadata, is considered safe since it reflects an editorial choice regarding what type of information each agent has access to.

\subsection{Implications for Interprofessional Education}

The risk of hallucination when using LLMs in medicine is particularly acute for clinical practice. While retrieval does not solve the problem of hallucinations, it does mitigate this risk because of the smaller surface area, as the evidence base used to feed the models is curated and validated. The partitioning of literature by agents provides a second layer of editing – the social work agent does not have access to the same body of literature as the physician assistant agent, which makes conflicts between agents more educational. It remains to be seen whether this is beneficial to learning outcomes, but this is something we cannot determine solely based on our metrics.

\subsection{Limitations}

Some important limitations must be addressed. \textit{Corpus size}: With only 116 data chunks,
this study’s reported statistics cannot be extended to a corpus of an order of magnitude more items. \textit{Query set}: This evaluation dataset is deliberately limited in scope to being an engineering benchmark for RAG system development; it is not meant to be a statement on generalization capabilities. The evaluation covers all ten topics of the opioid taxonomy and contains real-world student search queries but is too small to do statistical significance tests. Increasing the query size and acquiring relevance judgments from domain experts is essential. \textit{Scope of validation}: Response quality, pedagogic performance, or clinical accuracy were not measured, only the retrieval quality. \textit{Profession filter evaluation}: Reported results are obtained without profession
filtering and require relevance judgments per agent to properly evaluate the filtering effect. \textit{Single domain}: The data corpus consists of only opioid use disorder cases, and other clinical domains have not yet been tested.

\subsection{Threats to Validity}

The query for evaluation was formulated by the authors of the paper along with the reviewers, which creates an inherent risk of annotator bias in small corpus-based RAG evaluation. The ground-truth values for relevance were provided with the corpus and could be based on the same assumptions as the corpus itself. It would be useful to have a blind evaluation performed by an outside party. Generator GPT-5.2 is employed; performance may not carry over to smaller open-weight models.

%%=============================================================%%
\section{Conclusion and Future Work}\label{sec:conclusion}

The paper detailed our most recent improvements to BRIDGE, a retrieval-augmented, multi-agent interprofessional education system in healthcare. Apart from its RAG pipeline, BRIDGE utilizes a novel system for dynamic prompt engineering that includes prompt refinement, real-time RAG chunk injection, user profile personalization, and
referral-aware prompting to address limitations of static prompts in the baseline system. We also evaluated design choices that impact the response quality of the retrieval layer of the system. Five chunking techniques, four token budgets, four retrieval depths, seven retrievers, and two metadata filtering mechanisms were considered in our evaluation over an opioid education corpus. Recursive chunking $k=7$ for 7000 tokens provided the best balance between recall and prompt size. Among the retrievers, BGE-M3 and BM25 saturated MRR. The former scored the highest F1
and is therefore our preferred option, while the latter provides a good and computationally efficient backup. Topic filtering failed at present; profession filtering is still necessary, as it encodes an explicit editorial separation that fits into the multiple-perspective design of IPE.

The next steps include the following. \textit{Expansion of the evidence base through continuous ingestion}. The evidence base we considered is limited on purpose; further expanding it to include other specialties and scenarios and implementing a procedure for automatic tagging of professions within the corpus are the logical next steps. \textit{Clinical validation experiment}. A user study with healthcare students to evaluate whether the multi-agent RAG approach is more effective at teaching learners than the single agent or traditional IPE format is required. \textit{Multi-domain evaluation}. An extension of the evaluation to diabetes management and pediatric emergency care is planned in order to generalize design choices made during our experiment. \textit{Adaptive prompting.} Prompts could be made adaptive to learner responses in lieu of the static prompts used currently.

In conclusion, BRIDGE is a teaching tool, not a contribution to retrieval systems. The value of this work lies in evaluating the decisions involved in the development of an IPE system.

%%=============================================================%%
%% Backmatter
%%=============================================================%%
\backmatter

\bmhead{Acknowledgements}
The authors thank Dr. Andrew Wiss, supervising author and primary 
investigator for the BRIDGE project, of the Department of Health 
Policy and Management at The George Washington University. We also thank Dr. Amir Jafari of the Department of Data Science at the George Washington University for supporting this work. Inquiries regarding the broader BRIDGE research agenda may be directed to Dr. Andrew Wiss (awiss@gwu.edu).

\section*{Declarations}

\begin{itemize}
\item \textbf{Funding:} No funding was received for conducting this study.
\item \textbf{Conflict of interest:} The authors have no relevant financial or non-financial interests to disclose.
\item \textbf{Ethics approval and consent to participate:} This study was  conducted under Institutional Review Board (IRB) approval at The George Washington University (IRB\#: NCR256298 and NCR203171). 
Consent to participate was obtained from all participants.
\item \textbf{Consent for publication:} Not applicable.
\item \textbf{Data availability:} The opioid-education corpus and evaluation queries used in this study are available from the corresponding author upon reasonable request.
\item \textbf{Materials availability:} Not applicable.
\item \textbf{Code availability:} The RAG evaluation scripts are available from the corresponding author upon reasonable request.
\item \textbf{Author contribution:} All named authors contributed to the study conception and design. Corpus construction, extraction pipeline development, retrieval experiments, and statistical analysis were performed by Anusha Umashankar and Hari Prasannaa Thangavel Ravi. The first draft of the manuscript was written by Anusha Umashankar and Hari Prasannaa Thangavel Ravi. All commented on previous versions of the manuscript. All authors read and approved the final manuscript.
\end{itemize}

%%=============================================================%%
%% Appendix
%%=============================================================%%
\begin{appendices}
\section{Comprehensive Chunking Results}\label{appendix:chunking}

Table~\ref{tab:comprehensive_summary} reports every chunking-method, token-target, and retrieval-depth combination evaluated in Section~\ref{sec:results}, for 80 configurations in total.

\begin{sidewaystable}
\caption{Comprehensive chunking and retrieval summary across methods, token targets, and $k$ values. Latency in ms and build and index times in s.}
\label{tab:comprehensive_summary}

\rowcolors{2}{gray!15}{white}
\begin{tabular}{@{}lrrrrrrrrrr@{}}
\toprule
Method & Tokens & $k$ & Hit Rate & First Hit Rank & Latency & Chunks & Build & Index & Avg Tok & p95 Tok \\
\midrule
fixed & 3000 & 3 & 0.733 & 1.455 & 0.082 & 66 & 0.081 & 3.053 & 2684.0 & 3000.0 \\
fixed & 3000 & 5 & 0.733 & 1.455 & 0.079 & 66 & 0.081 & 3.053 & 2684.0 & 3000.0 \\
fixed & 3000 & 7 & 0.867 & 2.231 & 0.084 & 66 & 0.081 & 3.053 & 2684.0 & 3000.0 \\
fixed & 3000 & 9 & 0.867 & 2.231 & 0.078 & 66 & 0.081 & 3.053 & 2684.0 & 3000.0 \\
fixed\_overlap & 3000 & 3 & 0.533 & 1.625 & 0.092 & 76 & 0.144 & 3.991 & 2804.5 & 3000.0 \\
fixed\_overlap & 3000 & 5 & 0.533 & 1.625 & 0.115 & 76 & 0.144 & 3.991 & 2804.5 & 3000.0 \\
fixed\_overlap & 3000 & 7 & 0.533 & 1.625 & 0.110 & 76 & 0.144 & 3.991 & 2804.5 & 3000.0 \\
fixed\_overlap & 3000 & 9 & 0.600 & 2.444 & 0.115 & 76 & 0.144 & 3.991 & 2804.5 & 3000.0 \\
recursive & 3000 & 3 & 0.667 & 1.500 & 0.106 & 71 & 0.159 & 3.493 & 2495.3 & 2977.0 \\
recursive & 3000 & 5 & 0.733 & 1.818 & 0.120 & 71 & 0.159 & 3.493 & 2495.3 & 2977.0 \\
recursive & 3000 & 7 & 0.733 & 1.818 & 0.099 & 71 & 0.159 & 3.493 & 2495.3 & 2977.0 \\
recursive & 3000 & 9 & 0.800 & 2.333 & 0.102 & 71 & 0.159 & 3.493 & 2495.3 & 2977.0 \\
semantic & 3000 & 3 & 0.467 & 1.429 & 0.106 & 180 & 65.663 & 8.313 & 984.3 & 1127.7 \\
semantic & 3000 & 5 & 0.467 & 1.429 & 0.096 & 180 & 65.663 & 8.313 & 984.3 & 1127.7 \\
semantic & 3000 & 7 & 0.533 & 2.000 & 0.094 & 180 & 65.663 & 8.313 & 984.3 & 1127.7 \\
semantic & 3000 & 9 & 0.533 & 2.000 & 0.082 & 180 & 65.663 & 8.313 & 984.3 & 1127.7 \\
sentence\_pack & 3000 & 3 & 0.733 & 1.455 & 0.089 & 69 & 0.247 & 3.295 & 2567.6 & 2972.0 \\
sentence\_pack & 3000 & 5 & 0.800 & 1.750 & 0.078 & 69 & 0.247 & 3.295 & 2567.6 & 2972.0 \\
sentence\_pack & 3000 & 7 & 0.867 & 2.077 & 0.081 & 69 & 0.247 & 3.295 & 2567.6 & 2972.0 \\
sentence\_pack & 3000 & 9 & 0.867 & 2.077 & 0.068 & 69 & 0.247 & 3.295 & 2567.6 & 2972.0 \\
fixed & 5000 & 3 & 0.667 & 1.300 & 0.517 & 43 & 0.098 & 2.205 & 4119.6 & 5000.0 \\
fixed & 5000 & 5 & 0.733 & 1.545 & 0.065 & 43 & 0.098 & 2.205 & 4119.6 & 5000.0 \\
fixed & 5000 & 7 & 0.800 & 2.000 & 0.077 & 43 & 0.098 & 2.205 & 4119.6 & 5000.0 \\
fixed & 5000 & 9 & 0.867 & 2.462 & 0.072 & 43 & 0.098 & 2.205 & 4119.6 & 5000.0 \\
fixed\_overlap & 5000 & 3 & 0.800 & 1.500 & 0.090 & 49 & 0.143 & 2.480 & 4288.7 & 5000.0 \\
fixed\_overlap & 5000 & 5 & 0.800 & 1.500 & 0.092 & 49 & 0.143 & 2.480 & 4288.7 & 5000.0 \\
fixed\_overlap & 5000 & 7 & 0.800 & 1.500 & 0.098 & 49 & 0.143 & 2.480 & 4288.7 & 5000.0 \\
fixed\_overlap & 5000 & 9 & 0.867 & 2.077 & 0.097 & 49 & 0.143 & 2.480 & 4288.7 & 5000.0 \\
recursive & 5000 & 3 & 0.667 & 1.500 & 0.106 & 46 & 0.130 & 2.287 & 3851.1 & 4930.2 \\
recursive & 5000 & 5 & 0.667 & 1.500 & 0.086 & 46 & 0.130 & 2.287 & 3851.1 & 4930.2 \\
recursive & 5000 & 7 & 0.733 & 1.909 & 0.094 & 46 & 0.130 & 2.287 & 3851.1 & 4930.2 \\
recursive & 5000 & 9 & 0.867 & 3.000 & 0.076 & 46 & 0.130 & 2.287 & 3851.1 & 4930.2 \\
semantic & 5000 & 3 & 0.600 & 1.444 & 0.095 & 114 & 74.590 & 5.282 & 1554.3 & 1825.6 \\
semantic & 5000 & 5 & 0.667 & 1.800 & 0.062 & 114 & 74.590 & 5.282 & 1554.3 & 1825.6 \\
semantic & 5000 & 7 & 0.667 & 1.800 & 0.082 & 114 & 74.590 & 5.282 & 1554.3 & 1825.6 \\
semantic & 5000 & 9 & 0.667 & 1.800 & 0.081 & 114 & 74.590 & 5.282 & 1554.3 & 1825.6 \\
sentence\_pack & 5000 & 3 & 0.733 & 1.273 & 0.087 & 45 & 0.233 & 2.309 & 3936.8 & 4949.4 \\
sentence\_pack & 5000 & 5 & 0.733 & 1.273 & 0.081 & 45 & 0.233 & 2.309 & 3936.8 & 4949.4 \\
sentence\_pack & 5000 & 7 & 0.733 & 1.273 & 0.070 & 45 & 0.233 & 2.309 & 3936.8 & 4949.4 \\
sentence\_pack & 5000 & 9 & 0.867 & 2.385 & 0.081 & 45 & 0.233 & 2.309 & 3936.8 & 4949.4 \\

\botrule
\multicolumn{11}{r}{\textit{Continued on next page}} \\
\botrule
\end{tabular}

\end{sidewaystable}

\begin{sidewaystable}
\begin{tabular}{@{}lrrrrrrrrrr@{}}
\toprule
Method & Tokens & $k$ & Hit Rate & First Hit Rank & Latency & Chunks & Build & Index & Avg Tok & p95 Tok \\
\midrule

fixed & 7000 & 3 & 0.667 & 1.600 & 0.085 & 35 & 0.104 & 1.946 & 5061.3 & 7000.0 \\
fixed & 7000 & 5 & 0.733 & 1.818 & 0.080 & 35 & 0.104 & 1.946 & 5061.3 & 7000.0 \\
fixed & 7000 & 7 & 0.867 & 2.615 & 0.073 & 35 & 0.104 & 1.946 & 5061.3 & 7000.0 \\
fixed & 7000 & 9 & 0.933 & 3.000 & 0.080 & 35 & 0.104 & 1.946 & 5061.3 & 7000.0 \\
fixed\_overlap & 7000 & 3 & 0.733 & 1.455 & 0.096 & 38 & 0.134 & 2.117 & 5472.2 & 7000.0 \\
fixed\_overlap & 7000 & 5 & 0.800 & 1.667 & 0.105 & 38 & 0.134 & 2.117 & 5472.2 & 7000.0 \\
fixed\_overlap & 7000 & 7 & 0.867 & 2.000 & 0.101 & 38 & 0.134 & 2.117 & 5472.2 & 7000.0 \\
fixed\_overlap & 7000 & 9 & 0.933 & 2.429 & 0.095 & 38 & 0.134 & 2.117 & 5472.2 & 7000.0 \\
recursive & 7000 & 3 & 0.800 & 1.250 & 0.082 & 35 & 0.149 & 1.955 & 5061.4 & 6955.5 \\
recursive & 7000 & 5 & 0.800 & 1.250 & 0.089 & 35 & 0.149 & 1.955 & 5061.4 & 6955.5 \\
recursive & 7000 & 7 & 0.867 & 1.615 & 0.077 & 35 & 0.149 & 1.955 & 5061.4 & 6955.5 \\
recursive & 7000 & 9 & 0.867 & 1.615 & 0.095 & 35 & 0.149 & 1.955 & 5061.4 & 6955.5 \\
semantic & 7000 & 3 & 0.533 & 1.875 & 0.072 & 85 & 72.785 & 3.189 & 2084.4 & 2423.8 \\
semantic & 7000 & 5 & 0.600 & 2.111 & 0.096 & 85 & 72.785 & 3.189 & 2084.4 & 2423.8 \\
semantic & 7000 & 7 & 0.600 & 2.111 & 0.097 & 85 & 72.785 & 3.189 & 2084.4 & 2423.8 \\
semantic & 7000 & 9 & 0.667 & 2.800 & 0.090 & 85 & 72.785 & 3.189 & 2084.4 & 2423.8 \\
sentence\_pack & 7000 & 3 & 0.667 & 1.500 & 0.081 & 35 & 0.251 & 1.923 & 5061.6 & 6934.6 \\
sentence\_pack & 7000 & 5 & 0.733 & 1.727 & 0.081 & 35 & 0.251 & 1.923 & 5061.6 & 6934.6 \\
sentence\_pack & 7000 & 7 & 0.800 & 2.083 & 0.077 & 35 & 0.251 & 1.923 & 5061.6 & 6934.6 \\
sentence\_pack & 7000 & 9 & 0.867 & 2.615 & 0.073 & 35 & 0.251 & 1.923 & 5061.6 & 6934.6 \\
fixed & 10000 & 3 & 0.733 & 1.364 & 0.074 & 26 & 0.100 & 1.317 & 6813.2 & 10000.0 \\
fixed & 10000 & 5 & 0.800 & 1.667 & 0.081 & 26 & 0.100 & 1.317 & 6813.2 & 10000.0 \\
fixed & 10000 & 7 & 0.800 & 1.667 & 0.085 & 26 & 0.100 & 1.317 & 6813.2 & 10000.0 \\
fixed & 10000 & 9 & 0.933 & 2.643 & 0.079 & 26 & 0.100 & 1.317 & 6813.2 & 10000.0 \\
fixed\_overlap & 10000 & 3 & 0.733 & 1.273 & 0.085 & 28 & 0.179 & 1.562 & 7183.7 & 10000.0 \\
fixed\_overlap & 10000 & 5 & 0.733 & 1.273 & 0.098 & 28 & 0.179 & 1.562 & 7183.7 & 10000.0 \\
fixed\_overlap & 10000 & 7 & 0.867 & 2.077 & 0.097 & 28 & 0.179 & 1.562 & 7183.7 & 10000.0 \\
fixed\_overlap & 10000 & 9 & 0.933 & 2.500 & 0.090 & 28 & 0.179 & 1.562 & 7183.7 & 10000.0 \\
recursive & 10000 & 3 & 0.600 & 1.333 & 0.091 & 26 & 0.124 & 1.147 & 6813.4 & 9899.0 \\
recursive & 10000 & 5 & 0.800 & 2.167 & 0.092 & 26 & 0.124 & 1.147 & 6813.4 & 9899.0 \\
recursive & 10000 & 7 & 0.867 & 2.538 & 0.084 & 26 & 0.124 & 1.147 & 6813.4 & 9899.0 \\
recursive & 10000 & 9 & 0.933 & 2.929 & 0.092 & 26 & 0.124 & 1.147 & 6813.4 & 9899.0 \\
semantic & 10000 & 3 & 0.867 & 1.538 & 0.078 & 61 & 60.698 & 2.319 & 2904.3 & 3486.0 \\
semantic & 10000 & 5 & 0.933 & 1.786 & 0.091 & 61 & 60.698 & 2.319 & 2904.3 & 3486.0 \\
semantic & 10000 & 7 & 0.933 & 1.786 & 0.071 & 61 & 60.698 & 2.319 & 2904.3 & 3486.0 \\
semantic & 10000 & 9 & 0.933 & 1.786 & 0.077 & 61 & 60.698 & 2.319 & 2904.3 & 3486.0 \\
sentence\_pack & 10000 & 3 & 0.733 & 1.455 & 0.084 & 26 & 0.236 & 1.194 & 6813.6 & 9900.2 \\
sentence\_pack & 10000 & 5 & 0.800 & 1.667 & 0.075 & 26 & 0.236 & 1.194 & 6813.6 & 9900.2 \\
sentence\_pack & 10000 & 7 & 0.867 & 2.000 & 0.076 & 26 & 0.236 & 1.194 & 6813.6 & 9900.2 \\
sentence\_pack & 10000 & 9 & 0.933 & 2.500 & 0.090 & 26 & 0.236 & 1.194 & 6813.6 & 9900.2 \\

\botrule
\end{tabular}
\end{sidewaystable}

\end{appendices}

%%===========================================================================================%%
%% References
%%===========================================================================================%%
\clearpage
\begin{thebibliography}{99}

\bibitem{wiss2025groupthink}
Wiss, A., Showstark, M., Dobbeck, K., Pattershall-Geide, J., Zschaebitz, E., Joosten-Hagye, D., ... \& Embry, E. (2025). Utilizing Generative AI to Counter Learner Groupthink by Introducing Controversy in Collaborative Problem-Based Learning Settings. Online Learning Journal.

\bibitem{wiss2025aaipe}
Wiss, A., Joosten-Hagye, D., Pattershall-Geide, J., Showstark, M., Zschaebitz, E., Potter, K., ... \& Brooks, P. (2025). Development of the AI Acceptance Scale for Interprofessional Education (AAIPE) and collaborative practice settings. Journal of Interprofessional Education \& Practice, 40, 100752. 

\bibitem{reeves2016ipe}
Reeves S, Palaganas J, Zierler B (2017) An updated synthesis of review evidence of interprofessional education. J Allied Health 46(1):56--61

\bibitem{who2010ipe}
World Health Organization (2010) Framework for Action on Interprofessional Education and Collaborative Practice. WHO Press, Geneva

\bibitem{park2023generative}
Park JS, O'Brien JC, Cai CJ, Morris MR, Liang P, Bernstein MS (2023) Generative agents: interactive simulacra of human behavior. In: Proc. 36th Annual ACM Symposium on User Interface Software and Technology (UIST '23). \url{https://doi.org/10.1145/3586183.3606763}

\bibitem{wu2023autogen}
Wu Q, Bansal G, Zhang J, Wu Y, Li B, Zhu E, Jiang L, Zhang X, Zhang S, Liu J, Awadallah AH, White RW, Burger D, Wang C (2023) AutoGen: enabling next-gen LLM applications via multi-agent conversation. arXiv preprint arXiv:2308.08155

\bibitem{ji2023survey}
Ji Z, Lee N, Frieske R, Yu T, Su D, Xu Y, Ishii E, Bang YJ, Madotto A, Fung P (2023) Survey of hallucination in natural language generation. ACM Comput Surv 55(12):1--38. \url{https://doi.org/10.1145/3571730}

\bibitem{singhal2023medpalm}
Singhal K, Azizi S, Tu T, Mahdavi SS, Wei J, Chung HW, Scales N, Tanwani A, Cole-Lewis H, Pfohl S, Payne P, Seneviratne M, Gamble P, Kelly C, Babiker A, Sch\"{a}rli N, Chowdhery A, Mansfield P, Demner-Fushman D, Ag\"{u}era y Arcas B, Webster D, Corrado GS, Matias Y, Chou K, Gottweis J, Tomasev N, Liu Y, Rajkomar A, Barral J, Semturs C, Karthikesalingam A, Natarajan V (2023) Large language models encode clinical knowledge. Nature 620:172--180. \url{https://doi.org/10.1038/s41586-023-06291-2}

\bibitem{lewis2020rag}
Lewis P, Perez E, Piktus A, Petroni F, Karpukhin V, Goyal N, K\"{u}ttler H, Lewis M, Yih W, Rockt\"{a}schel T, Riedel S, Kiela D (2020) Retrieval-augmented generation for knowledge-intensive NLP tasks. In: Advances in Neural Information Processing Systems 33 (NeurIPS 2020)

\bibitem{gao2023ragsurvey}
Gao Y, Xiong Y, Gao X, Jia K, Pan J, Bi Y, Dai Y, Sun J, Wang H (2023). Retrieval-augmented generation for large language models: a survey. arXiv preprint arXiv:2312.10997

\bibitem{izacard2021fid}
Izacard G, Grave E (2021) Leveraging passage retrieval with generative models for open domain question answering. In: Proc. 16th Conference of the European Chapter of the Association for Computational Linguistics (EACL)

\bibitem{langchain2023}
LangChain (2023) Text splitters documentation. \url{https://python.langchain.com/docs/modules/data_connection/document_transformers/}. Accessed 10 Feb 2026

\bibitem{kamradt2023}
Kamradt G (2023) Semantic chunking for retrieval-augmented generation. Technical report, Greg Kamradt AI Research

\bibitem{gunther2024late}
G\"{u}nther M, Milliken L, Geuter J, Mastrapas G, Wang B, Xiao H (2024) Late chunking: contextual chunk embeddings using long-context embedding models. arXiv preprint arXiv:2409.04701

\bibitem{reimers2019sbert}
Reimers N, Gurevych I (2019) Sentence-BERT: sentence embeddings using Siamese BERT-networks. In: Proc. 2019 Conference on Empirical Methods in Natural Language Processing (EMNLP)

\bibitem{song2020mpnet}
Song K, Tan X, Qin T, Lu J, Liu T-Y (2020) MPNet: masked and permuted pre-training for language understanding. In: Advances in Neural Information Processing Systems 33 (NeurIPS 2020)

\bibitem{xiao2023bge}
Xiao S, Liu Z, Zhang P, Muennighoff N (2023) C-pack: packaged resources to advance general Chinese embedding. arXiv preprint arXiv:2309.07597

\bibitem{chen2024bgem3}
Chen J, Xiao S, Zhang P, Luo K, Lian D, Liu Z (2024) BGE M3-embedding: multi-lingual, multi-functionality, multi-granularity text embeddings through self-knowledge distillation. arXiv preprint arXiv:2402.03216

\bibitem{wang2022e5}
Wang L, Yang N, Huang X, Jiao B, Yang L, Jiang D, Majumder R, Wei F (2022) Text embeddings by weakly-supervised contrastive pre-training. arXiv preprint arXiv:2212.03533

\bibitem{robertson2009bm25}
Robertson S, Zaragoza H (2009) The probabilistic relevance framework: BM25 and beyond. Found Trends Inf Retr 3(4):333--389. \url{https://doi.org/10.1561/1500000019}

\bibitem{thakur2021beir}
Thakur N, Reimers N, R\"{u}ckl\'{e} A, Srivastava A, Gurevych I (2021) BEIR: a heterogeneous benchmark for zero-shot evaluation of information retrieval models. In: Proc. NeurIPS 2021 Track on Datasets and Benchmarks

\bibitem{li2023camel}
Li G, Hammoud H, Itani H, Khizbullin D, Ghanem B (2023) CAMEL: communicative agents for "mind" exploration of large language model society. In: Advances in Neural Information Processing Systems 36 (NeurIPS 2023)

\bibitem{lee2020biobert}
Lee J, Yoon W, Kim S, Kim D, Kim S, So CH, Kang J (2020) BioBERT: a pre-trained biomedical language representation model for biomedical text mining. Bioinformatics 36(4):1234--1240. \url{https://doi.org/10.1093/bioinformatics/btz682}

\bibitem{stevens2006virtual}
Stevens A, Hernandez J, Johnsen K, Dickerson R, Raij A, Harrison C, DiPietro M, Allen B, Ferdig R, Foti S, Jackson J, Shin M, Cendan J, Watson R, Duerson M, Lok B, Cohen M, Wagner P, Lind DS (2006) The use of virtual patients to teach medical students history taking and communication skills. Am J Surg 191(6):806--811

\end{thebibliography}

\end{document}